% ═══════════════════════════════════════════════════════════════
% ML-08a: Musterlösung - Viajar -- Reisen und Verkehrsmittel
% Abgeleitet aus LE-08a (Abschnitt 9: Lösungen)
% ═══════════════════════════════════════════════════════════════

\documentclass[11pt, a4paper]{article}

\usepackage[utf8]{inputenc}
\usepackage[T1]{fontenc}
\usepackage[ngerman]{babel}
\usepackage{helvet}
\renewcommand{\familydefault}{\sfdefault}

\usepackage[margin=2cm]{geometry}
\usepackage{enumitem}
\usepackage{tabularx}
\usepackage{booktabs}

\usepackage{xcolor}
\usepackage{tcolorbox}

\usepackage{fancyhdr}
\usepackage{lastpage}

% Farben
\definecolor{spanisch}{HTML}{c41e3a}
\definecolor{grau}{HTML}{888888}
\definecolor{hellgrau}{HTML}{f5f5f5}
\definecolor{orange}{HTML}{b35c00}
\definecolor{loesung}{HTML}{c62828}

\newcommand{\fachfarbe}{spanisch}

% Variablen
\newcommand{\thema}{Viajar -- Reisen und Verkehrsmittel}
\newcommand{\sequenz}{08}
\newcommand{\block}{a}
\newcommand{\fach}{Spanisch}
\newcommand{\klasse}{7}
\newcommand{\maxpunkte}{24}

% Kopf-/Fußzeile
\pagestyle{fancy}
\fancyhf{}
\renewcommand{\headrulewidth}{0.4pt}
\renewcommand{\footrulewidth}{0.4pt}
\lhead{\fach{} -- Klasse \klasse}
\chead{\textcolor{loesung}{\textbf{ML-\sequenz\block{} (MUSTERLÖSUNG)}}}
\rhead{}
\lfoot{}
\rfoot{Seite \thepage\ von \pageref{LastPage}}

% Befehle
\newcommand{\punktebox}[1]{\hfill\fbox{\small \_/#1}}
\newcommand{\loesung}[1]{\textcolor{loesung}{\textbf{#1}}}
\newcommand{\luecke}[1]{\rule{#1}{0.4pt}}
\newcommand{\es}[1]{\textcolor{\fachfarbe}{\textbf{#1}}}

% Merkkasten
\newtcolorbox{merkkasten}[1][]{
    colback=hellgrau,
    colframe=\fachfarbe,
    boxrule=1pt,
    arc=0pt,
    left=8pt, right=8pt, top=8pt, bottom=8pt,
    fonttitle=\bfseries\color{\fachfarbe},
    title=#1
}

% Hinweis
\newtcolorbox{hinweis}{
    colback=white,
    colframe=orange,
    boxrule=1pt,
    arc=0pt,
    left=5pt, right=5pt, top=5pt, bottom=5pt,
    fonttitle=\bfseries,
    title=Hinweis
}

% Aufgabe
\newenvironment{aufgabe}[1]{%
    \vspace{10pt}
    \noindent\textbf{\textcolor{\fachfarbe}{Aufgabe #1}}
    \vspace{5pt}
    \begin{list}{}{\leftmargin=0pt}
    \item[]
}{%
    \end{list}
}

% ═══════════════════════════════════════════════════════════════
\begin{document}

% Kopfbereich
\begin{center}
    {\Large\bfseries\textcolor{\fachfarbe}{Arbeitsblatt: \thema}}\\[0.3em]
    {\large\textcolor{loesung}{\textbf{--- MUSTERLÖSUNG ---}}}
\end{center}

\vspace{5pt}

\noindent
\begin{tabularx}{\textwidth}{|X|X|X|}
\hline
\textbf{Name:} \textcolor{loesung}{Musterschüler/in} &
\textbf{Klasse:} \klasse &
\textbf{Datum:} \textcolor{loesung}{XX.XX.XXXX} \\
\hline
\end{tabularx}

\vspace{15pt}

% ═══════════════════════════════════════════════════════════════
% MERKKASTEN 1 (mit Lösungen)
% ═══════════════════════════════════════════════════════════════

\begin{merkkasten}[Merkkasten: Los medios de transporte (Verkehrsmittel)]

\begin{center}
\begin{tabular}{|l|l|l|}
\hline
\textbf{Spanisch} & \textbf{Deutsch} & \textbf{Beispiel} \\
\hline
\loesung{el avión} & das Flugzeug & Vamos \loesung{en} avión. \\
\hline
\loesung{el tren} & der Zug & El tren sale a las ocho. \\
\hline
\loesung{el coche} & das Auto & Viajamos en coche. \\
\hline
\loesung{el autobús} & der Bus & Tomo el autobús. \\
\hline
\loesung{el barco} & das Schiff & El barco va a las islas. \\
\hline
\end{tabular}
\end{center}

\end{merkkasten}

\vspace{10pt}

% ═══════════════════════════════════════════════════════════════
% MERKKASTEN 2 (mit Lösungen)
% ═══════════════════════════════════════════════════════════════

\begin{merkkasten}[Merkkasten: ir a + Infinitiv (nahe Zukunft)]

\begin{center}
\begin{tabular}{|l|l|l|}
\hline
\textbf{Person} & \textbf{ir} & \textbf{+ a + Infinitiv} \\
\hline
yo & \loesung{voy} & a nadar \\
\hline
tú & \loesung{vas} & a viajar \\
\hline
él/ella & \loesung{va} & a hacer senderismo \\
\hline
nosotros & \loesung{vamos} & a tomar el sol \\
\hline
\end{tabular}
\end{center}

\vspace{0.5em}
\textbf{Bedeutung:} \textit{Voy a nadar} = Ich werde schwimmen / Ich gehe schwimmen.
\end{merkkasten}

\vspace{15pt}

% ═══════════════════════════════════════════════════════════════
% AUFGABE 1 (mit Lösungen)
% ═══════════════════════════════════════════════════════════════

\begin{aufgabe}{1}
Schreibe die spanischen Wörter für die Verkehrsmittel. \punktebox{5}
\end{aufgabe}

\begin{enumerate}[label=\alph*)]
    \item Flugzeug: \loesung{el avión}
    \item Zug: \loesung{el tren}
    \item Auto: \loesung{el coche}
    \item Bus: \loesung{el autobús}
    \item Schiff: \loesung{el barco}
\end{enumerate}

\textit{\textcolor{grau}{Je 1P pro richtige Antwort}}

\vspace{15pt}

% ═══════════════════════════════════════════════════════════════
% AUFGABE 2 (mit Lösungen)
% ═══════════════════════════════════════════════════════════════

\begin{aufgabe}{2}
Verbinde die Orte mit den passenden Aktivitäten. \punktebox{4}
\end{aufgabe}

\begin{center}
\begin{tabular}{rcl}
la playa & \loesung{------} & \loesung{nadar, tomar el sol} \\[0.8em]
la montaña & \loesung{------} & \loesung{hacer senderismo} \\[0.8em]
el hotel & \loesung{------} & \loesung{descansar} \\[0.8em]
el camping & \loesung{------} & \loesung{hacer una barbacoa} \\
\end{tabular}
\end{center}

\textit{\textcolor{grau}{Je 1P pro richtige Zuordnung}}

\vspace{15pt}

% ═══════════════════════════════════════════════════════════════
% AUFGABE 3 (mit Lösungen)
% ═══════════════════════════════════════════════════════════════

\begin{aufgabe}{3}
Ergänze die Sätze mit der richtigen Form von \textit{ir a + Infinitiv}. \punktebox{6}
\end{aufgabe}

\begin{enumerate}[label=\alph*)]
    \item Yo \loesung{voy a nadar} en la playa. \textit{(schwimmen)}
    \item Nosotros \loesung{vamos a viajar} en avión. \textit{(reisen)}
    \item Tú \loesung{vas a hacer} senderismo. \textit{(machen)}
    \item Ellos \loesung{van a tomar} el sol. \textit{(nehmen/sich sonnen)}
    \item Ella \loesung{va a ir} a la montaña. \textit{(gehen)}
    \item Vosotros \loesung{vais a quedaros} en el hotel. \textit{(bleiben/quedarse)}
\end{enumerate}

\textit{\textcolor{grau}{Je 1P pro richtige Form (0,5P ir + 0,5P Infinitiv)}}

\vspace{15pt}

% ═══════════════════════════════════════════════════════════════
% AUFGABE 4 (mit Lösungen)
% ═══════════════════════════════════════════════════════════════

\begin{aufgabe}{4}
Vervollständige den Dialog. Nutze die Redemittel. \punktebox{6}
\end{aufgabe}

\noindent
\textbf{A:} \es{¿Adónde vas a ir de vacaciones?}

\textbf{B:} \loesung{Voy a ir a la playa. / Voy a ir a España.}

\textbf{A:} \es{¿Cómo vas a viajar?}

\textbf{B:} \loesung{Voy a viajar en avión. / Voy a viajar en coche.}

\textbf{A:} \loesung{¿Qué vas a hacer?}

\textbf{B:} \loesung{Voy a nadar y a tomar el sol.}

\vspace{0.5em}
\textit{\textcolor{grau}{Je 1-2P pro sinnvolle Ergänzung, max. 6P}}

\vspace{15pt}

% ═══════════════════════════════════════════════════════════════
% AUFGABE 5 (Beispiellösung)
% ═══════════════════════════════════════════════════════════════

\begin{aufgabe}{5}
Schreibe einen kurzen Reiseplan (mind. 5 Sätze). \punktebox{3}
\end{aufgabe}

\textbf{Beispiellösung:}

\loesung{Este verano voy a ir a España.}

\loesung{Voy a viajar en avión con mi familia.}

\loesung{Vamos a quedarnos en un hotel cerca de la playa.}

\loesung{Voy a nadar todos los días.}

\loesung{También voy a tomar el sol y a comer tapas.}

\vspace{0.5em}
\textit{\textcolor{grau}{Bewertung: Ziel (0,5P), Verkehrsmittel (0,5P), Unterkunft (0,5P), Aktivitäten (1P), Sprachliche Korrektheit (0,5P)}}

\vspace{20pt}
\hrule
\vspace{10pt}

\begin{flushright}
\textbf{Punkte:} \loesung{\maxpunkte} / \maxpunkte
\end{flushright}

\end{document}
